% Source: https://artofproblemsolving.com/community/c3h2472293_march_2021_problem_set
% Stable identity: aops-topic-2472293-q07
Setting $y=0$ gives $p(x)+p(0)+1=p(x)$, so $p(0)=-1$. Setting $y=1$ gives
\[p(x+1)=p(x)+x+2\qquad\text{for all real }x.\]
Let $g(k)=\dfrac{k^2+3k-2}{2}$. Then $g(0)=-1$ and $g(k+1)-g(k)=\dfrac{2k+1+3}{2}=k+2$, so $g$ satisfies the same recursion. Starting from $p(0)=g(0)$, induction upwards (using $p(k+1)=p(k)+k+2$) and downwards (using $p(k)=p(k+1)-k-2$) gives $p(k)=g(k)$ for every integer $k$.

Such a function exists: $p(x)=\dfrac{x^2+3x-2}2$ satisfies $p(1)=1$ and
\[p(x)+p(y)+1=\dfrac{x^2+y^2+3x+3y-2}2=\dfrac{(x+y)^2+3(x+y)-2}2-xy=p(x+y)-xy.\]
Now for an integer $k$,
\[p(k)=k\iff k^2+3k-2=2k\iff k^2+k-2=0\iff(k+2)(k-1)=0.\]
So $k=1$ or $k=-2$, and exactly $\boxed{1}$ integer $k\ne1$ (namely $k=-2$) has $p(k)=k$.
